% Part: propositional-logic
% Chapter: propositional-logic 
% Section: completeness

\documentclass[../../../include/open-logic-section]{subfiles}

\begin{document}

\olfileid[hi]{pl}{prp}{com}

\olsection{प्रतिज्ञप्ति तर्कशास्त्र की पूर्णता}

\begin{defn}
\ollabel{def:MaxCon}
!!{formula}s का समुच्चय $\Gamma$ \emph{अधिकतम संगत} है यदि वह संगत है और
यदि $\Delta$ ऐसा संगत समुच्चय है कि $\Gamma \subseteq \Delta$, तो $\Gamma
= \Delta$।
\end{defn}

\begin{lem}[सत्यता लेम्मा]
  \ollabel{lem:truth}
मान लें $\Gamma$ अधिकतम संगत है; तब:
\begin{enumerate}
\item \ollabel{lem:truth-1} $!A \in \Gamma$ तभी और केवल तभी जब $\lnot !A
  \notin \Gamma$;
\item \ollabel{lem:truth-2} $\Gamma \Proves !A$ तभी और केवल तभी जब $!A \in \Gamma$; 
\item \ollabel{lem:truth-3} $!A \lif !B \in \Gamma$ तभी और केवल तभी जब
  या तो $!A \notin \Gamma$ या $!B \in \Gamma$।
\end{enumerate}
\end{lem}

\begin{proof}
\begin{enumerate}
\item मान लें $!A \in \Gamma$। यदि $\lnot !A \in \Gamma$ भी, तो $\Gamma$
  असंगत है; और यदि न $!A$ न $\lnot!A$, $\Gamma$ में है, तो
  \olref[axd]{prop:phi} से $\Gamma\cup\{!A\}$ और $\Gamma \cup \{\lnot
  !A\}$ में से कोई एक संगत है, जिसका अर्थ है कि $\Gamma$ अधिकतम नहीं है।
  दूसरी दिशा भी इसी प्रकार सिद्ध होती है।
\item यदि $!A \in \Gamma$, तो निस्संदेह $\Gamma \Proves !A$। दूसरी ओर
  मान लें $\Gamma \Proves !A$। यदि $!A \notin \Gamma$, तो
  \olref{lem:truth-1} से $\lnot !A \in \Gamma$। अतः $\Gamma \Proves !A$
  और $\Gamma \Proves \lnot !A$, जो विरोधाभास है।
\item अभ्यास।
\end{enumerate}
\end{proof}

\olref{lem:truth}\olref{lem:truth-2} की बाएँ-से-दाएँ दिशा दिखाती है कि
$\Gamma$ \emph{निगमनात्मक रूप से संवृत} है, अर्थात् सभी~$!A$ के लिए यदि
$\Gamma \Proves !A$, तो $!A \in \Gamma$।

\begin{prob}
  \olref{lem:truth} का प्रमाण पूरा कीजिए।
\end{prob}

\begin{thm}[पूर्णता]
\ollabel{thm:completeness} 
यदि $\Gamma$ संगत है, तो $\Gamma$ संतुष्टनीय है।
\end{thm}

\begin{proof}
मान लें $!A_0$, $!A_1$,~\dots भाषा के सभी !!{formula}s की सर्वग्राही सूची
है। हम !!{formula}s के समुच्चयों का बढ़ता अनुक्रम $\Gamma_0$,
$\Gamma_1$,~\dots निम्न प्रकार पुनरावर्ती रूप से परिभाषित करते हैं:
\begin{align*}
\Gamma_0 & = \Gamma\\
\Gamma_{n+1} &  =
\begin{cases}
  \Gamma_n \cup \{!A_n\} & \text{यदि $\Gamma_n \cup \{!A_n\}$ संगत है;}\\
  \Gamma_n \cup \{\lnot !A_n\} & \text{अन्यथा।}
\end{cases}
\end{align*}
फिर परिभाषित करें:
\[
\Gamma^* = \bigcup_{0\le n}\Gamma_n.
\]
अब प्रमाण बारी-बारी से निम्नलिखित तथ्य स्थापित करके आगे बढ़ता है:
\begin{enumerate}
\item प्रत्येक $n$ के लिए समुच्चय $\Gamma_n$ संगत है ($n$ पर आगमन और
  \olref[axd]{prop:phi} के उपयोग से);
\item $\Gamma^*$ संगत है;
\item $\Gamma^*$ अधिकतम है।
\end{enumerate}
फिर मूल्यांकन $v$ को $v(p_i) = \True$ तभी और केवल तभी रखकर परिभाषित करें जब
$p_i \in \Gamma^*$। इसके बाद $!A$ पर आगमन से दिखाया जाता है कि अधिक जटिल
!!{sentence}s के लिए भी $\Gamma^*$ की सदस्यता, $v$ के अनुसार सत्यता से मेल
खाती है: $\overline{v}(!A) = \True$ तभी और केवल तभी जब $!A \in \Gamma^*$।
विशेषतः $v$, $\Gamma^*$ को संतुष्ट करता है और चूँकि $\Gamma \subseteq
\Gamma^*$, इसलिए इच्छानुसार $\Gamma$ भी संतुष्टनीय है।
\end{proof}

\begin{cor}
यदि $\Gamma \Entails !A$, तो $\Gamma \Proves !A$।
\end{cor}

\begin{proof}
यदि $\Gamma\Proves/ !A$, तो \olref[axd]{prop:prov-incons} से $\Gamma \cup
\{\lnot !A\}$ संगत है। अतः पूर्णता प्रमेय से $\Gamma \cup \{\lnot !A\}$
संतुष्टनीय है और $\Gamma \Entails/ !A$।
\end{proof}

\begin{prop}[संहतता प्रमेय]
\ollabel{prop:compactness}
$\Gamma$ तभी और केवल तभी संतुष्टनीय है जब $\Gamma$ का प्रत्येक \emph{परिमित}
उपसमुच्चय $\Gamma_0$ संतुष्टनीय हो।
\end{prop}

\begin{proof}
$\Gamma$ तभी और केवल तभी असंतुष्टनीय है जब वह असंगत है, तभी और केवल तभी जब
  $\Gamma$ का कोई परिमित उपसमुच्चय $\Gamma_0$ असंगत है, और तभी और केवल तभी
  जब $\Gamma$ का कोई परिमित उपसमुच्चय $\Gamma_0$ असंतुष्टनीय है।
\end{proof}

\begin{cor}
$\Gamma \Entails !A$ तभी और केवल तभी जब $\Gamma$ के किसी परिमित उपसमुच्चय
$\Gamma_0$ के लिए $\Gamma_0 \Entails !A$।
\end{cor}

\end{document}
